\documentclass[11pt,a4paper]{article}

% --- Page margins and paragraphs ---
\usepackage[margin=1in]{geometry}
\usepackage{parskip}

% --- Basic math and typesetting ---
\usepackage{amsmath,amssymb,amsthm}
\usepackage{mathtools} % provides \coloneqq
\usepackage{url}
\usepackage{xcolor} %

% --- Theorem and definition environments ---
\theoremstyle{definition}
\newtheorem{definition}{Definition}
\newtheorem{example}{Example}

\theoremstyle{plain}
\newtheorem{theorem}{Theorem}
\newtheorem{proposition}{Proposition}
\newtheorem{lemma}{Lemma}

\begin{document}

% ==================================================
% Title and author information
% ==================================================
\begin{center}
    {\LARGE \textbf{Mathematics Method for Physics--Lecture 11}} \\[0.6em]
    {\large Computing Residues and the Residue at Infinity}
\end{center}

\vspace{-0.5em}
\noindent\rule{\textwidth}{0.5pt}
\vspace{0.3em}

\begin{center}
    Instructor: Zhaolong Han, SOCE, Shanghai Jiao Tong University;
\end{center}

\vspace{-0.6em}
\noindent Email: han.arkey@sjtu.edu.cn

\vspace{0.4em}
\noindent\rule{\textwidth}{0.5pt}

\vspace{0.8em}
% Blue hint/explanation text
\noindent { Lecture notes for November 8.}

\vspace{1.5em}

% ==================================================
% Main content
% ==================================================
\section{From the residue theorem to computation}

Lecture 10 showed that a contour integral can be evaluated from the coefficients of $(z-z_0)^{-1}$ at the enclosed singularities. This lecture explains how to obtain those coefficients without expanding an entire Laurent series.

We use $\mathrm{i}^2=-1$ and counterclockwise orientation for a positive finite contour. Every local contour below encloses only the singularity under discussion. Worked exercises and the original practice calculations are retained in the accompanying exercises files.

\section{Computing a residue from a Laurent series}

If $f$ is holomorphic in $0<|z-z_0|<R$, write

\[
f(z)=\sum_{k=-\infty}^{\infty}C_k(z-z_0)^k,\qquad
\operatorname{Res}[f,z_0]\coloneqq C_{-1}.
\]

One needs only the terms that can contribute to exponent $-1$. For a product of convergent Laurent series on a common annulus, a term of degree $p$ in one factor combines with a term of degree $q$ in the other when $p+q=-1$. Absolute convergence on a fixed interior circle justifies collecting those terms.

This approach applies to essential singularities as well as poles. At a removable singularity, there are no negative powers and the residue is zero.

\section{Residue formulas at poles}

\subsection{A pole of order n}

\textbf{Theorem.} Suppose $z_0$ is a pole of order $n\ge1$. Define the holomorphic extension

\[
h(z)\coloneqq(z-z_0)^nf(z).
\]

Then

\[
\operatorname{Res}[f,z_0]
=\frac{h^{(n-1)}(z_0)}{(n-1)!}
=\frac1{(n-1)!}\lim_{z\to z_0}
\frac{\mathrm{d}^{\,n-1}}{\mathrm{d}z^{n-1}}
\left[(z-z_0)^nf(z)\right].
\]

\textbf{Proof.} The Laurent expansion and the corresponding Taylor expansion of $h$ are

\[
\begin{aligned}
f(z)&=\frac{C_{-n}}{(z-z_0)^n}+\cdots+\frac{C_{-1}}{z-z_0}
+C_0+C_1(z-z_0)+\cdots,\\
h(z)&=C_{-n}+\cdots+C_{-1}(z-z_0)^{n-1}
+C_0(z-z_0)^n+\cdots.
\end{aligned}
\]

Differentiate the second line $n-1$ times. Lower-degree terms disappear, and every term of degree at least $n$ still has a factor $z-z_0$. Thus

\[
h^{(n-1)}(z)=C_{-1}(n-1)!+O(z-z_0).
\]

Letting $z\to z_0$ gives the formula. The differentiation acts on the entire product $(z-z_0)^nf(z)$.

\subsection{Simple and double poles}

For $n=1$ and $n=2$, respectively, the theorem becomes

\[
\begin{aligned}
\operatorname{Res}[f,z_0]&=\lim_{z\to z_0}(z-z_0)f(z),\\
\operatorname{Res}[f,z_0]&=\lim_{z\to z_0}
\frac{\mathrm{d}}{\mathrm{d}z}\left[(z-z_0)^2f(z)\right].
\end{aligned}
\]

An equivalent criterion for a simple pole is that $(z-z_0)f(z)$ extends holomorphically and has a nonzero value at $z_0$. Its value is then the residue.

\subsection{A quotient with a simple zero in the denominator}

Let $f=p/q$, where $p$ and $q$ are holomorphic near $z_0$, $q(z_0)=0$, and $q'(z_0)\ne0$. Taylor expansion gives

\[
q(z)=q'(z_0)(z-z_0)+O((z-z_0)^2).
\]

Substitute into the simple-pole limit to obtain

\[
\operatorname{Res}[p/q,z_0]=\frac{p(z_0)}{q'(z_0)}.
\]

When $p(z_0)\ne0$, the singularity is a simple pole. When $p(z_0)=0$, cancellation makes it removable and the same expression gives zero. This criterion requires a simple zero of $q$.

\subsection{Choosing a calculation method}

Identify the local order before choosing a pole formula. Cancellation in a numerator can reduce the order suggested by the denominator. For an essential singularity, use a Laurent expansion or a coefficient identity. For a quotient with a simple denominator zero, the derivative formula above usually avoids unnecessary expansion.

\section{Residues at infinity}

\subsection{Exterior expansions and orientation}

Assume $f$ is holomorphic for $|z|>R$. Its Laurent expansion on this exterior annulus is

\[
f(z)=\sum_{n=1}^{\infty}\frac{C_{-n}}{z^n}+C_0+\sum_{n=1}^{\infty}C_nz^n.
\]

Let $C_\rho^+$ be the counterclockwise circle $|z|=\rho>R$, and let $C_\rho^-$ have the opposite orientation. The coefficient calculation from Lecture 10 gives

\[
\oint_{C_\rho^+}f(z)\,\mathrm{d}z=2\pi\mathrm{i}C_{-1},
\qquad
\oint_{C_\rho^-}f(z)\,\mathrm{d}z=-2\pi\mathrm{i}C_{-1}.
\]

\textbf{Definition.} The residue at infinity is

\[
\operatorname{Res}[f,\infty]\coloneqq
\frac1{2\pi\mathrm{i}}\oint_{C_\rho^-}f(z)\,\mathrm{d}z
=-C_{-1}.
\]

The coefficient is taken from the exterior expansion. In particular,

\[
\oint_{C_\rho^+}f(z)\,\mathrm{d}z=-2\pi\mathrm{i}\operatorname{Res}[f,\infty].
\]

A large clockwise circle is positively oriented around infinity in its local coordinate. Under $\zeta=1/z$, it maps to a small counterclockwise circle around zero. This accounts for the orientation convention.

\subsection{Transforming the differential}

The local coordinate at infinity is $\zeta=1/z$. Since $z=1/\zeta$, its differential satisfies

\[
\mathrm{d}z=-\frac1{\zeta^2}\,\mathrm{d}\zeta,
\qquad
f(z)\,\mathrm{d}z=-\frac1{\zeta^2}f(1/\zeta)\,\mathrm{d}\zeta.
\]

Thus residues transform together with the differential. Applying the finite-point residue definition in the $\zeta$ coordinate gives

\[
\operatorname{Res}[f,\infty]
=-\operatorname{Res}\left[\zeta^{-2}f(1/\zeta),0\right].
\]

To check the sign and exponent directly, substitute into the exterior Laurent expansion:

\[
\begin{aligned}
f(1/\zeta)&=\cdots+C_{-2}\zeta^2+C_{-1}\zeta+C_0+C_1/\zeta+\cdots,\\
\zeta^{-2}f(1/\zeta)&=\cdots+C_{-2}+C_{-1}/\zeta+C_0/\zeta^2+\cdots.
\end{aligned}
\]

The coefficient of $\zeta^{-1}$ in the second line is $C_{-1}$. The minus sign comes from $\mathrm{d}z=-\zeta^{-2}\mathrm{d}\zeta$.

\section{The sum of all residues}

\textbf{Theorem.} Suppose $f$ is holomorphic on $\mathbb C$ except at finitely many isolated singularities $z_1,\ldots,z_N$. Then infinity is isolated and

\[
\sum_{k=1}^{N}\operatorname{Res}[f,z_k]+\operatorname{Res}[f,\infty]=0.
\]

\textbf{Proof.} Choose a counterclockwise circle $C_\rho^+$ containing all finite singularities. The residue theorem and the definition at infinity evaluate the same contour integral:

\[
\frac1{2\pi\mathrm{i}}\oint_{C_\rho^+}f(z)\,\mathrm{d}z
=\sum_{k=1}^{N}\operatorname{Res}[f,z_k]
=-\operatorname{Res}[f,\infty].
\]

Rearrangement gives the claim. Essential singularities are allowed as long as they are isolated. The finite set hypothesis ensures that a circle can enclose all finite singularities.

This identity is especially helpful when the exterior expansion is simpler than the individual local expansions. It gives the sum of the finite residues directly.

\section{A limit formula and its hypotheses}

\textbf{Proposition.} If $f$ is holomorphic in an exterior region and $\lim_{z\to\infty}zf(z)=A\in\mathbb C$, then

\[
\operatorname{Res}[f,\infty]=-A.
\]

\textbf{Proof.} Set $H(\zeta)\coloneqq \zeta^{-1}f(1/\zeta)$. The assumed limit means that $H$ has a removable singularity at zero, with Taylor expansion

\[
H(\zeta)=A+b_1\zeta+b_2\zeta^2+\cdots.
\]

Consequently,

\[
f(1/\zeta)=A\zeta+b_1\zeta^2+\cdots,
\qquad
f(z)=\frac A z+\frac{b_1}{z^2}+\cdots.
\]

Hence $C_{-1}=A$, proving the result. This proof also shows that $f(z)\to0$ under the stated hypothesis.

The weaker assumption $f(z)\to L$ only establishes removability at infinity. It does not determine the coefficient of $1/z$. For example, $L+A/z$ has the same finite limit $L$ for every $A$, while its residue at infinity equals $-A$.

\section{Further reading}

For the contour interpretation and finite-point residue theorem, see Mathews and Howell's open textbook. Complex Analysis, Section 8.1\footnote{\url{https://complexanalysis.org/web/sec_residue-thm.html}}


\end{document}
